\documentclass[11pt]{article}
\usepackage[a4paper,margin=28mm]{geometry}
\usepackage{amsmath,amssymb}
\usepackage{booktabs}
\usepackage{enumitem}
\usepackage[hidelinks]{hyperref}
\usepackage[T1]{fontenc}

\newcommand{\EA}{(1+1)\text{-EA}}
\newcommand{\runtimeSup}{\operatorname{runtimeSup}}

\title{Formal Runtime Analysis of Evolutionary Algorithms in Lean\\
\large Current Verified State and Development Frontier}
\author{ECRuntime project}
\date{25 September 2026}

\begin{document}
\maketitle

\noindent\textbf{Frozen at Stage~79.} This note is not extended after the 547-module
merge. The current paper-by-paper coverage is \texttt{docs/current/coverage.md}.

\begin{abstract}
This note records the current mathematical and formalization state of the
\texttt{ECRuntime} Lean project after the clean Stage~79 fix~1 verification. The project
provides reusable stochastic-process and runtime-analysis machinery for evolutionary
computation and applies it to OneMax, LeadingOnes, and arbitrary positive ordered linear
functions. Stage~79 fix~1 is the latest compiler-promoted checkpoint: it promotes the
corrected factorial-moment chain through the closed gamma-tail identity, the corrected
suffix-ratio geometry, pointwise nonnegative accepted left contribution, exact singleton
left gain, and event-weighted lower bound $(p_1/n)g_i$. There is no active Stage~80
theorem candidate in this repository-state checkpoint.
\end{abstract}

\section{Verification boundary}

The project is pinned to Lean~4.32.1 and Mathlib~v4.32.1. The returned Stage~79 fix~1 verification is warning/error-free for \texttt{StaticUnaryUnbiasedCorrectedLeftGain}, the complete root build, repository export, and source-hygiene checks. Consequently Stage~79 fix~1 is compiler promoted.

Consequently, for arbitrary positive nonincreasing linear weights, the finite Witt line
is compiler promoted through
\[
  0\le c\le\log 3,\qquad n\ge4,
\]
with final runtime interfaces on $0<c\le\log3$.

For
\[
  q_n(c)=\frac{c}{n}\left(1-\frac{c}{n}\right)^{n-1},
\]
the fixed-scale asymptotic layer proves for every fixed $c>0$
\[
  n q_n(c)\longrightarrow c e^{-c},
  \qquad
  (nq_n(c))^{-1}\longrightarrow \frac{e^c}{c}.
\]
The leading factor $e^c/c$ is minimized at $c=1$.

\section{Reusable formal infrastructure}

\subsection{Trajectory semantics}

The repository contains finite path weights and consistency, stochastic transition
kernels, Ionescu--Tulcea infinite trajectories, finite-marginal identification, and
least-index first-hitting-time semantics. These bridges allow finite-state recurrence
arguments and sample-path runtime statements to share a common formal foundation.

\subsection{Runtime theory}

Reusable analysis includes additive and multiplicative drift, fitness-level bounds,
stochastic domination, geometric waiting variables and finite sums, level-exit/residence
infrastructure, and geometric/shifted-exponential runtime tails.

\section{Completed application: LeadingOnes}

The LeadingOnes development is exposed through
\texttt{ECRuntime.LeadingOnesMainResults}. It includes first-crossing symmetry, exact
entry distributions, residence and occupation identities, renewal, and exact expected
runtime with scaled-rate and standard-rate interfaces.

\section{Completed finite-boundary Witt line}

For $p=c/n$, the closed finite C3 condition is
\[
  \left(1-\frac cn\right)^{-(n-1)}+\frac cn\le3.
\]
Its limiting boundary is $e^c\le3$. Stage~66 proved the exact endpoint $c=\log3$ by a
dimension-uniform logarithmic argument, and Stage~67 propagated the resulting certificate
through the generic expectation and tail machinery.

\texttt{WittMutationC3Envelope.lean} owns endpoint-independent C3 monotonicity;
\texttt{WittMutationCertifiedRuntime.lean} owns the C3-based finite expectation and tail
consequences; and \texttt{WittMutationLogThreeRuntime.lean} is the final thin
specialization. Historical rational endpoints beyond the retained $21/20$ compatibility
layer have been removed from the maintained proof graph.

This branch is now frozen unless a later theorem family exposes a reusable abstraction
that genuinely improves it.

\section{New theorem family: static unary-unbiased mutation}

The next target follows the static unary-unbiased linear-function setting. A unary
unbiased mutation operator on $\{0,1\}^n$ is characterized by choosing a Hamming radius
and then choosing uniformly among masks of that radius. For runtime analysis the central
parameters are
\[
  p_1=\Pr(\text{flip exactly one bit})
\]
and the expected mutation strength
\[
  \chi=\mathbb E[\text{number of flipped bits}].
\]

\subsection{Stage 68 fix 1 promoted operator interface}

Stage~68 fix~1 compiler-promotes \texttt{StaticUnaryUnbiasedMutation.lean}. It represents
a static unary-unbiased operator by a normalized nonnegative law on mutation masks that
is invariant under mask cardinality. The module derives the radius masses, their
normalization, the one-bit mass $p_1$, the mean radius $\chi$, the exact specified
singleton mass $p_1/n$, aggregate single-zero-flip mass, the independent-bit embedding,
and the induced elitist stochastic transition.

\subsection{Stage 69 fix 1 promoted: singleton linear drift}

Stage~69 fix~1 compiler-promotes \texttt{StaticUnaryUnbiasedLinearDrift.lean}. It deliberately formalizes
only the distribution-generic part of the linear-function progress argument:
\begin{itemize}[nosep]
  \item a mask-sum formula for one-step drift under the static unary-unbiased transition;
  \item the named per-coordinate singleton rate $p_1/n$;
  \item exact contribution $(p_1/n)V(x)$ of zero-singleton events for arbitrary
        linear-error coefficients;
  \item specialization to the bounded rank potential;
  \item an exact decomposition of rank drift into the singleton term and one
        non-singleton remainder;
  \item a coefficient-sensitive objective-error multiplicative-drift baseline at rate
        $p_1/n$.
\end{itemize}

The candidate does not yet assert a coefficient-independent runtime theorem for general
static unary-unbiased mutation. That requires a lower bound on the non-singleton
remainder.

\section{Why this is the next efficient step}

The 2024 static unary-unbiased result has asymptotic scale
\[
  (1\pm o(1))\frac{n\log n}{p_1}
\]
under additional conditions, including control of the expected mutation strength
$\chi$. The repository already contains most of the generic trajectory/runtime machinery
and extensive linear-potential infrastructure. The missing bridge is therefore not a new
probability space but an operator abstraction exposing $p_1$ and the multi-bit radius
profile cleanly.

Stage~68 created precisely that interface. Stage~69 now consumes it in the exact singleton branch while keeping all multi-bit corrections explicit.

\section{Stage 70 fix 1 promoted: coordinate marginals}

Stage~70 fix~1 compiler-promotes \texttt{StaticUnaryUnbiasedFlipMarginals.lean}. Its primary theorem is
\[
  \Pr(i\text{ flips})=\frac{\chi}{n}
\]
for every specified coordinate in positive dimension. The module also promotes the unconditional shifted radius moment and the second factorial radius moment
\[
  \sum_k k(k-1)p_k,
\]
which is the correct distributional quantity for two-coordinate flip events.

\section{Stage 71 fix 1 promoted: pair marginals}

Stage~71 fix~1 compiler-promotes \texttt{StaticUnaryUnbiasedPairMarginals.lean}. It defines the event mass that two specified coordinates flip, proves equality of this mass for all distinct coordinate pairs by swap symmetry, bounds it by the promoted coordinate marginal $\chi/n$, and rewrites the second factorial radius moment as the mask expectation of $|M|(|M|-1)$.

\section{Stage 72 fix 1 promoted: exact pair normalization}

Stage~72 fix~1 compiler-promotes \texttt{StaticUnaryUnbiasedPairNormalization.lean}. It double-counts ordered distinct coordinate pairs inside each mask, identifies their aggregate mass with $\mathbb E[K(K-1)]$, and combines this identity with Stage~71 exchangeability to obtain
\[
  \Pr(i,j\text{ flip})=\frac{\mathbb E[K(K-1)]}{n(n-1)}
\]
for distinct $i,j$ and $n\ge2$. When $\chi>0$, dividing by the promoted coordinate marginal gives the exact size-biased conditional ratio
\[
  \frac{\Pr(i,j\text{ flip})}{\Pr(j\text{ flips})}
  =\frac{\mathbb E[K(K-1)]}{\chi(n-1)}.
\]
No bound replacing the factorial moment by a function of $\chi$ alone is asserted.

The first Stage~72 compiler run exposed only local proof-engineering issues. Stage~72 fix~1 repaired those points, and the focused module plus complete root build then verified warning/error-free. The exact pair normalization and size-biased ratio are therefore compiler promoted.

\section{Stage 73 promoted: generic event-mass bridge}

Stage~73 compiler-promotes \texttt{StaticUnaryUnbiasedAcceptedEventBridge.lean}. For a finite mask event $A$ that contains the singleton $\{i\}$ and is supported on masks flipping $i$, the event mass of outcomes in $A$ that also flip a distinct coordinate $j$ is at most the promoted pair mass, while $A$ itself has mass at least $p_1/n$. Hence, for $p_1>0$,
\[
  \Pr(j\text{ flips}\mid A)\le
  \frac{\mathbb E[K(K-1)]}{p_1(n-1)}.
\]
The theorem is kept at the finite event-mass level. Stage~74 instantiates it with the concrete accepted leftmost-error pivot events of the linear-function proof without introducing an incorrect mean-only conditioning law.


\section{Stage 74 fix 1 promoted: concrete accepted pivot events}

Stage~74 fix~1 compiler-promotes \texttt{StaticUnaryUnbiasedLinearAcceptedEvents.lean}. Because this repository maximizes positive linear functions toward the all-one string and numbers the most important coordinate first, the paper's leftmost flipped error bit is represented as a currently-zero pivot $i$ such that no smaller-index zero-to-one flip occurs. The event additionally requires elitist acceptance.

The promoted module proves pivot support, pivot zero status, acceptance, the leftmost/minimal-index property, singleton inclusion for positive linear weights, and pairwise disjointness of distinct pivot events. Under $p_1>0$ it proves positive event mass and instantiates the Stage~73 bridge to obtain
\[
  \Pr(j\text{ flips}\mid A_i)\le
  \frac{\mathbb E[K(K-1)]}{p_1(n-1)}.
\]

\section{Stage 75 fix 1 promoted: accepted-event left/right loss decomposition}

Stage~75 fix~1 compiler-promotes \texttt{StaticUnaryUnbiasedLinearLossDecomposition.lean}. It partitions currently-one coordinates around an abstract anchor and proves the exact accepted weighted-error decomposition
\[
  \Delta=\Delta_L+\Delta_R.
\]
For nonnegative coefficients it also sums the promoted Stage~74 per-coordinate estimate over the right region, retaining the exact factorial-moment coefficient
\[
  \frac{\mathbb E[K(K-1)]}{p_1(n-1)}.
\]

\section{Stage 76 fix 1 promoted: corrected potential calibration}

Stage~76 fix~1 compiler-promotes \texttt{StaticUnaryUnbiasedCorrectedPotential.lean}. Combining the promoted event-mass bound $\chi/n$ with the Stage~75 conditional coefficient and normalizing by the singleton rate $p_1/n$ yields the corrected dimensionless penalty
\[
  \rho=\frac{\chi\,\mathbb E[K(K-1)]}{(n-1)p_1^2}.
\]
The promoted module proves this calibration identity exactly, defines the geometric baseline $1+\alpha\rho$, and in the repository's reversed coordinate convention sets
\[
  \gamma_i=(1+\alpha\rho)^{n-1-i}.
\]
For positive nonincreasing objective weights it constructs the clipped coefficient through the finite suffix minimum
\[
  g_i=\min_{j\ge i}\gamma_j\frac{w_i}{w_j},
\]
retaining a concrete minimizing suffix coordinate as the reversed anchor. It proves anchor membership/minimality, $1\le g_i\le\gamma_i$, nonnegativity, and the all-ones target characterization of the resulting potential. The first Stage~76 build exposed only elaboration/linter issues; Fix~1 made suffix finiteness explicit and removed a redundant tactic, after which the focused and full builds verified warning/error-free.

\section{Stage 77 fix 1 and Stage 78 fix 2 promoted: closed corrected gamma tail}

Stage~77 fix~1 compiler-promotes \texttt{StaticUnaryUnbiasedCorrectedPotentialEnvelope.lean}. It proves the exact anchor-candidate representation
\[
  g_i=\gamma_{a(i)}\frac{w_i}{w_{a(i)}}
\]
and, for positive nonincreasing objective weights, the lower envelope $\gamma_{a(i)}\le g_i$. It also promotes the corrected-potential $L/R$ specialization and the reduction of the complete conditional right-region loss to the pure gamma suffix tail.

Stage~78 adds \texttt{StaticUnaryUnbiasedCorrectedGammaTail.lean}. With $\rho$ denoting the promoted corrected penalty ratio and $\gamma_i=(1+\alpha\rho)^{n-1-i}$, it evaluates the strict suffix geometrically and proves
\[
  \Gamma_{>k}(\alpha\rho)=\gamma_k-1.
\]
For $\alpha>0$ this implies $\rho\Gamma_{>a(i)}\le g_i/\alpha$. Combining this with the promoted calibration identity converts the Stage~77 conditional right-loss envelope into singleton-rate units. The next boundary is the actual pivot-event mass/left-gain assembly and then the global multiplicative-drift theorem.

The first Stage~78 pinned build found no mathematical counterexample: it failed only in the inverse finite-suffix reindexing constructor, where an `omega` call could not see through the newly constructed `Fin` value. Stage~78 fix~1 moved the natural-number side conditions outside that constructor. The next pinned build confirmed the constructor repair and localized the only remaining errors to the `left_inv` and `right_inv` laws of the same equivalence, where `omega` could not normalize the opaque composed maps. Stage~78 fix~2 exposes the two involution equalities explicitly before arithmetic automation. All Stage~78 theorem statements remain unchanged. The Stage~78-fix~2 focused and root builds then verified warning/error-free, so this closed gamma-tail layer is now promoted.

\section{Lemma 9 audit before further formalization}

The published 2024 proof bounds a conditional pair-flip probability in Lemma~9. Its displayed equation~(27), however, conditions on one specified coordinate being flipped and then averages $(k-1)/(n-1)$ against the original radius probabilities $p_k$. Conditioning on a specified flipped coordinate size-biases the radius law by a factor $k/\chi$, so this displayed step is not valid for an arbitrary radius law as written.

A simple diagnostic family shows why mean control alone cannot justify that intermediate inequality: take $p_1=1/2$, $p_n=1/(2n)$ and $p_0=1/2-1/(2n)$. Then $\chi=1$, but for two distinct coordinates the conditional probability that the first flips given that the second flips is $1/2$, whereas $\chi/(n-1)=1/(n-1)$. For $n\ge4$ the claimed generic conditional bound fails.

Therefore the repository does not encode the mean-only pair bound as a theorem. Stage~71 isolates pair-event symmetry, Stage~72 compiler-promotes exact pair normalization, and Stage~73 compiler-promotes the generic corrected event-mass coefficient based on $\mathbb E[K(K-1)]$. Stage~74 fix~1 now compiler-promotes the objective-specific accepted leftmost-error events and instantiates this bridge. Stage~75 fix~1 compiler-promotes the generic left/right decomposition and propagates the same factorial-moment coefficient over the whole right region. Stage~76 fix~1 now compiler-promotes the potential calibration directly to that verified coefficient, and Stage~77 reduces the right-loss term to a pure geometric tail. If later drift assembly does not recover a stronger estimate, the runtime theorem must retain a second-factorial-moment parameter or state an explicit additional moment hypothesis.

\section{Stage 79 fix 1 promoted: corrected left gain}

Stage~79 fix~1 compiler-promotes \texttt{StaticUnaryUnbiasedCorrectedLeftGain.lean}. The new layer proves
that the corrected potential-to-objective ratio is nondecreasing along a pivot suffix,
\[
  \frac{g_i}{w_i}\le \frac{g_q}{w_q}\qquad(i\le q),
\]
and that every coordinate on the left-loss side of the corrected anchor satisfies
\[
  \frac{g_j}{w_j}\le \frac{g_i}{w_i}\qquad(j\le a(i)).
\]
Together with objective acceptance these ratio envelopes imply the pointwise inequality
$\Delta_L(i)\ge0$ for every mask in the accepted pivot event.  The singleton pivot mask
has exactly $\Delta_L(i)=g_i$, so the event-weighted left contribution is at least
$(p_1/n)g_i$.  The first pinned Stage~79 build exposed only local proof-engineering
issues: an unused simplifier hypothesis, a hidden proof of $n>0$, and an elaboration
failure in the singleton lower-bound summation step. Stage~79 fix~1 repaired only those
points and preserved the complete theorem surface. The focused Stage~79-fix~1 build and
the complete root build then verified warning/error-free, so this corrected-left-gain
layer is now promoted. The remaining local task is to uncondition the promoted right-loss
proxy into the actual event-weighted right contribution and combine the two terms into
the corrected per-pivot drift inequality.


\section{Repository literature checkpoint after Stage 79 fix 1}

No Stage~80 theorem is introduced in this checkpoint.  The certified theorem inventory,
its mapping to the originating EC literature, the current limitations of each paper-level
formalization, and a refreshed survey of more than twenty reachable next papers are
recorded in
\texttt{docs/FORMALIZED\_RESULTS\_AND\_LITERATURE\_FRONTIER.md}.
The survey is dated 2026-09-25 and separates compiler-certified paper-level results,
certified supporting chains whose headline theorem is not yet closed, and generic
infrastructure.  In particular, it records explicitly that the 2024 static unary-unbiased
headline theorem has not yet been certified under the paper's published assumptions;
the repository continues with the exact factorial-moment correction exposed by the
Stage~70--79 chain.

\end{document}
